% ----------------------
\subsection{Section 129 (9 February 1843)}
% ----------------------
	\label{DC129}
	
	% -----
	\subsubsection{The original source for DC 129}
	% -----
	
		Note that the JSP has two original sources for the remarks that make up section 129. The first, which seems to be the source the section actually comes from, is \hyperref[EMS-1843-JS-9-FEB-1843-WC]{William Clayton's account}. The other source is \hyperref[EMS-1843-JS-9-FEB-1843-WR]{Willard Richards' account}.
	
	% -----
	\subsubsection{Historical Background}
	% -----
		\label{DC129-HB}
		
		% --
		\paragraph{JSP}
		% --
		
			``On 9 February 1843, JS delivered an instruction to guests he was entertaining at his home in Nauvoo, Illinois. JS apparently gave the remarks during a casual conversation with apostles Orson Pratt and Parley P. Pratt and others who had called on him in the afternoon. Though they were not delivered as a formal discourse, his teachings were significant enough that William Clayton and Willard Richards, who were both present on this occasion, wrote them down. Prefacing his remarks, JS `related some of his history' and then instructed those in attendance how they `might know whether any administration was from God.'

			``JS's remarks touched on a variety of doctrinal teachings that he had shared with other groups at various times since June 1839. Some of his comments were similar to teachings regarding the discerning of spirits that he originally shared with some members of the Quorum of the Twelve Apostles on 27 June 1839. On that occasion, JS spoke of a key to `detect Satan when he transforms himself nigh unto an angel of light.' He explained that a person might distinguish between an angel of God and Satan by trying to shake hands. Still imprisoned in Missouri in June 1839, Parley P. Pratt missed the meeting and JS's instruction. Orson Pratt was also likely absent when JS gave the instruction, having arrived in Columbia, Missouri, by 1 July 1839 to assist with Parley's escape from prison. Shortly after Parley escaped from custody, he and Orson arrived in Nauvoo in July 1839. The two brothers left for a mission to England shortly thereafter. Although Orson had returned to Nauvoo by early August 1841, he became disaffected from JS and the church in early summer 1842 and remained so until January 1843. Parley remained in England longer than the rest of the apostles and did not return to Nauvoo until 7 February 1843. This February 1843 meeting, therefore, may have been the first time JS shared these 1839 teachings with either of the Pratts. 

			``JS's remarks in February 1843 centered on an explanation of the ways to distinguish between various types of heavenly messengers and the devil. Among eighteenthand nineteenth-century Protestants, there were several competing ideas about the identification of angels, most of them based on a passage in the epistle to the Hebrews that mentions `an innumerable company of angels' and the `spirits of just men made perfect.' Some religionists held that these verses refer to `those translated to heaven in their bodies, and those raised from the dead after Christ's resurrection.' Others believed they refer more generally to `all in every age and nation who have feared God and wrought righteousness.' Theologians also debated the substance and materiality of angels. Eighteenth-century Swedish mystic Emmanuel Swedenborg suggested that angels were corporeal beings who had lived on the earth and could converse with men face-to-face. While those who believed angels were translated or risen beings seemingly believed that all angels were corporeal beings, some concluded that `angels have no corporeal forms.'314 Theologian Charles Buck explained that `as to the \textit{nature} of these beings we are told that they are spirits,' with the `more general opinion' being that `they are substances entirely spiritual.' At the same time, Buck allowed that `they can at any time assume bodies and appear in human shape.' 

			``JS suggested a new idea, which classified heavenly messengers as either resurrected corporeal beings or disembodied spirits awaiting resurrection. These distinctions may have appeared in Latter-day Saint theology as early as 1829, when the Book of Mormon suggested a difference between `angels and ministering spirits.' In JS's `new translation' of the Bible, however, Hebrews 1:7 explained that `angels are ministering spirits.' In his 27 June 1839 discourse, JS emphasized that `an angel of God (which is an angel of light) is a Saint with his resurrected body' but also noted that it was possible to be visited by deceased believers who were not yet resurrected. JS used this occasion in February 1843 to refine that explanation by distinguishing between `resurrected personages' and `the spirits of just men made perfect' who were still awaiting resurrection. 

			``In addition to teaching periodically about the ways to distinguish between types of angels, JS had demonstrated a long-standing interest in recognizing the differences between true and false spirits. A circa 8 March 1831 revelation urged the Saints to `beware lest ye are deceived' and to do `all things with prayer \& thanksgiving that ye may not be seduced by evil spirits or doctrines of Devils or the commandments of men for some are of men \& others of Devils.' A revelation the following May was even more specific, explaining, `There are many spirits which are false spirits which have gone forth in the Earth deceiving the world \& also Satan hath sought to deceive you that he might overthrow you.' In April 1842 JS reiterated this message in a lengthy editorial in the church newspaper, urging the Saints to `try the spirits.' Then, recounting portions of his own personal history to the Saints in a September 1842 letter, JS alluded to an early experience `on the banks of the Susquehanna' when the devil had appeared to him `as an Angel of light.'''
			
		% --
		\paragraph{HDDC}
		% --
		
			HDDC 1701--1704 has more background information and quotes more sources, but I'm not sure it goes beyond what's in the JSP.
					
	% -----
	\subsubsection{Versions}
	% -----
		\label{DC129-V}
		
		See the \href{https://drive.google.com/file/d/1goAvNz8jS4R8slYfMG_db55Ty2z_vmgK/view?usp=sharing}{sections comparison} document.
	
		\begin{compactdesc}
			\item[JSP] \hyperref[EMS-1843-JS-9-FEB-1843]{9 February 1843}
			\item[BC] ---
			\item[DC 1835] ---
			\item[DC 1844] ---
			\item[DN] 6:7 (23 April 1856), 49
			\item[MS] 20:33 (14 August 1858), 519
		\end{compactdesc}

	% -----
	\subsubsection{Section 129:1} % *DC129-1 
	% -----
		\label{DC129-1}

		THERE are two kinds of beings in heaven, namely: Angels, who are resurrected personages, having bodies of flesh and bones---

		% \begin{framed}
		% \scriptsize
			% \begin{compactdesc}
				% \item[JSP] 
				% % \item[BC] 
				% % \item[]
			% \end{compactdesc}
		% \end{framed}

	% -----
	\subsubsection{Section 129:2} % *DC129-2 
	% -----
		\label{DC129-2}

		For instance, Jesus said: \textit{Handle me and see, for a spirit hath not flesh and bones, as ye see me have.}

		% \begin{framed}
		% \scriptsize
			% \begin{compactdesc}
				% \item[JSP] 
				% % \item[BC] 
				% % \item[]
			% \end{compactdesc}
		% \end{framed}

	% -----
	\subsubsection{Section 129:3} % *DC129-3 
	% -----
		\label{DC129-3}

		Secondly: the spirits of just men made perfect, they who are not resurrected, but inherit the same glory.

		% \begin{framed}
		% \scriptsize
			% \begin{compactdesc}
				% \item[JSP] 
				% % \item[BC] 
				% % \item[]
			% \end{compactdesc}
		% \end{framed}

	% -----
	\subsubsection{Section 129:4} % *DC129-4 
	% -----
		\label{DC129-4}

		When a messenger comes saying he has a message from God, offer him your hand and request him to shake hands with you.

		% \begin{framed}
		% \scriptsize
			% \begin{compactdesc}
				% \item[JSP] 
				% % \item[BC] 
				% % \item[]
			% \end{compactdesc}
		% \end{framed}

	% -----
	\subsubsection{Section 129:5} % *DC129-5 
	% -----
		\label{DC129-5}

		If he be an angel he will do so, and you will feel his hand.

		% \begin{framed}
		% \scriptsize
			% \begin{compactdesc}
				% \item[JSP] 
				% % \item[BC] 
				% % \item[]
			% \end{compactdesc}
		% \end{framed}

	% -----
	\subsubsection{Section 129:6} % *DC129-6 
	% -----
		\label{DC129-6}

		If he be the spirit of a just man made perfect he will come in his glory; for that is the only way he can appear---

		% \begin{framed}
		% \scriptsize
			% \begin{compactdesc}
				% \item[JSP] 
				% % \item[BC] 
				% % \item[]
			% \end{compactdesc}
		% \end{framed}

	% -----
	\subsubsection{Section 129:7} % *DC129-7 
	% -----
		\label{DC129-7}

		Ask him to shake hands with you, but he will not move, because it is contrary to the order of heaven for a just man to deceive; but he will still deliver his message.

		% \begin{framed}
		% \scriptsize
			% \begin{compactdesc}
				% \item[JSP] 
				% % \item[BC] 
				% % \item[]
			% \end{compactdesc}
		% \end{framed}

	% -----
	\subsubsection{Section 129:8} % *DC129-8 
	% -----
		\label{DC129-8}

		If it be the devil as an angel of light, when you ask him to shake hands he will offer you his hand, and you will not feel anything; you may therefore detect him.

		% \begin{framed}
		% \scriptsize
			% \begin{compactdesc}
				% \item[JSP] 
				% % \item[BC] 
				% % \item[]
			% \end{compactdesc}
		% \end{framed}

	% -----
	\subsubsection{Section 129:9} % *DC129-9 
	% -----
		\label{DC129-9}

		These are three grand keys whereby you may know whether any administration is from God.

		% \begin{framed}
		% \scriptsize
			% \begin{compactdesc}
				% \item[JSP] 
				% % \item[BC] 
				% % \item[]
			% \end{compactdesc}
		% \end{framed}